\section{System}
\label{sec:system}

\subsection{Architecture overview}

The system is a deterministic skill chain that processes each 5-minute telemetry window. Figure~\ref{fig:architecture} (TBD) shows the data flow. A \emph{skill} is a typed function that reads from and writes to a shared \texttt{AgentContext}; skills are composable, deterministic, and individually testable. The default skill chain is fixed at compile time --- we deliberately do not use a learned planner because reviewer-reproducibility matters more than the marginal accuracy lift an LLM planner might offer.

The chain has eleven steps:

\begin{enumerate}
    \item \texttt{entity\_extract} --- parse window evidence to pull $(\text{service}, \text{component}, \text{error\_class}, \text{severity})$ tuples.
    \item \texttt{component\_filter} --- cheap pre-filter on the Jira memory, keeping only tickets sharing at least one entity with the window.
    \item \texttt{lexical\_similarity} --- BM25 scoring over the filtered candidate set.
    \item \texttt{log\_signature\_similarity} --- Move-A: replace the bi-encoder query and document text with a \emph{characteristic log line} per side. The characteristic line is extracted via term-frequency analysis over the raw log dump, yielding the highest-information line that a human engineer would also pick. This step aligns query and document vocabulary on both sides and is responsible for a +9\% Recall@5 lift over BM25-only in our experiments.
    \item \texttt{cross\_encoder\_rerank} --- a sentence-transformers MS-MARCO MiniLM cross-encoder reranks the BM25 top-20 by joint $(\text{query}, \text{doc})$ scoring. The cross-encoder's output is blended with the bi-encoder score at weight $0.6$ to preserve broader recall; we found that full replacement (weight $1.0$) lifted Hit@1 by 30\% relative but cost 18\% on Hit@5.
    \item \texttt{graph\_score} --- bridge-weighted scoring over an entity graph built from the corpus. Each candidate ticket's score gets a bonus for every entity shared with the window, weighted by how diagnostic that entity kind has been historically.
    \item \texttt{numeric\_blend} --- a gradient-boosted classifier fitted on the 94 derived numeric features (CPU\%, latency p99, error counts, k8s restart counts, etc.) emits a per-window triage probability. This signal is blended into the per-candidate combined score.
    \item \texttt{triage\_decide} --- final blend: $\text{triage\_score} = w_{\text{num}} \cdot \text{numeric\_score} + (1 - w_{\text{num}}) \cdot \max_{c} \text{combined}_c$, where $w_{\text{num}} = 0.80$ is tuned on the validation split.
    \item \texttt{novelty\_check} --- flag the incident as novel if $\max_c \text{combined}_c < 0.15$. This threshold is brittle; we discuss the resulting cold-start failure in Section~\ref{sec:limitations}.
    \item \texttt{graph\_traverse\_explain} --- produce a human-readable explanation citing the graph entities responsible for the top match.
    \item Final output: $\{\text{triage\_score}, \text{top-5 candidates}, \text{is\_novel}, \text{explanation}\}$.
\end{enumerate}

\subsection{The humanized Jira memory corpus (V2)}
\label{sec:corpus}

The retrieval-side memory is a corpus of 347 engineer-voice incident tickets, generated by a deterministic humanizer pipeline anchored to \textsc{TAWOS}'s empirical distributions of ticket length (mean 2{,}396 characters), comment count (mean 4.2 comments per ticket), and code-block ratio (90.2\% of tickets contain at least one fenced code block). Each ticket has:

\begin{itemize}
    \item A leading \texttt{description\_code} block --- engineer-vocabulary log lines from when the incident was filed. Placed first in the memory text so BM25 and short-document embeddings weight it.
    \item A timeline of persona-tagged steps (on-call, sub-system owner, SRE-lead), each step containing prose followed by a \texttt{body\_code} block of log lines the persona pasted.
    \item Standard incident metadata: affected services, components, severity, fault compatibility class.
\end{itemize}

The humanizer is sanitizer-verified: zero lab-terminology leaks (no occurrences of \texttt{chaos}, \texttt{synthetic}, \texttt{scenario}, \texttt{near-miss}, etc.) reach the indexed memory text. We measured this with a hard-coded regex sweep before every corpus build; any leak fails the build.

\subsection{Cross-encoder fine-tuning}

Starting from \texttt{cross-encoder/ms-marco-MiniLM-L-6-v2}, we fine-tune on 12{,}641 $(\text{query}, \text{doc}, \text{label})$ pairs built from the train+val splits:

\begin{itemize}
    \item \textbf{8{,}855 positives:} every gold-matching $(\text{window}, \text{ticket})$ pair extracted from the per-window matchings file. A single window can have multiple gold matches; each becomes its own positive pair.
    \item \textbf{3{,}786 hard negatives:} for each window, we retrieve the top-20 BM25 candidates excluding gold, then sample 3 per window as hard negatives. These are tickets that look superficially relevant under bag-of-words scoring but are not in the ground-truth set.
\end{itemize}

We train for 3 epochs with batch size 16, learning rate $2 \times 10^{-5}$, warmup ratio 0.1, and BCE loss with label smoothing disabled. Best checkpoint selected on validation BCE loss: epoch 1 (val loss 1.063, AP 0.737, F1 0.802). Epochs 2 and 3 overfit (val loss rises to 4.0).
